\subsubsection{Example: Load}
% TODO: remind that the repr is non-opt now
Consider the symbolic state $s_8$ before the execution of~\cref{line:llvm-load}:
\[
\begin{array}{c}
s_8 \triangleq \symstate
    {\location_9}
    {\{ ..., \nameid{gep2} \mapsto e \}}
    {[]}
    {[\mo]}
    {n < \bvconst{32}{2}} \\
e \triangleq \bvadd{\bvsort{64}}{\bvconst{64}{\addr}}{\bvzext{64}{\bvextract{32}{\bvzext{64}{n}}}} \\
\mo \triangleq \mobj{\bvval{\addr}{64}}{\bvval{2}{32}}{\arrayvalue}, \ \arrayvalue \triangleq \arraystore{\arraystore{\arrayconst{\arrayinitvalue}}{\bvconst{32}{0}}{\bvconst{8}{7}}}{\bvconst{32}{1}}{\arrayinitvalue}
\end{array}
\]
According to the symbolic semantics,
if $\mech{s}_9$ is derived from $s_8$ then there exist a model $m$ and a memory object $\mo' \triangleq \mobj{\addr'}{\size'}{\arrayvalue'}$ such that
$m \models n < \bvconst{32}{2}$, $\mo' \in \symmemory$, and $\symisinbounds{e}{\bvsort{8}}{\mo'}{m}$, and then:
% TODO: should add 0_{32}?
\[
\mech{s}_9 \triangleq \symstate
    {\location_{10}}
    {\{ ..., \nameid{gep2} \mapsto e, \nameid{v} \mapsto \arrayselect{\arrayvalue'}{\bvextract{32}{\bvsub{\bvsort{64}}{e}{\bvconst{64}{\addr}}}} \}}
    {[]}
    {[\mo]}
    {n < \bvconst{32}{2} \land \rangecond(e, 1, \mo')}
\]
Before dereferencing $e$, the \SE engine checks if $e$ may point to other memory objects,
by checking the satisfiability of: $n < \bvconst{32}{2} \land \neg \rangecond(e, 1, \mo)$.
Here, this constraint is unsatisfiable, so $e$ points only to $\mo$.
In particular, this implies that $\mo' = \mo$, and $\arrayvalue' = \arrayvalue$.
The symbolic state computed by the \SE engine is:
\[
s_9 \triangleq \symstate
    {\location_{10}}
    {\{ ..., \nameid{gep2} \mapsto e, \nameid{v} \mapsto \arrayselect{\arrayvalue}{\bvextract{32}{\bvsub{\bvsort{64}}{e}{\bvconst{64}{\addr}}}} \}}
    {[]}
    {[\mo]}
    {n < \bvconst{32}{2}}
\]
Now, we have to prove that the subtree of $s_8$ is safe.
We already know that $n < \bvconst{32}{2} \land \neg \rangecond(e, 1, \mo)$ is unsatisfiable,
so by~\Cref{def:sym-error-state}, $s_8$ is not an error state.
To prove that $\mech{s}_9 \equiv s_9$, it suffices to prove the equivalence of their path constraints.
Since $n < \bvconst{32}{2} \land \neg \rangecond(e, 1, \mo)$ is unsatisfiable,
this implies that $n < \bvconst{32}{2} \land \rangecond(e, 1, \mo) \equiv n < \bvconst{32}{2}$.

To handle the case of $\tload$ in general, we define the following interface lemma:
\begin{interfacelemma}
\label{lemma:interface-load}
Let $s \triangleq \symstate{\location}{\symlocalstore}{\symstack}{\symmemory}{\pc}$,
where \textup{$\instrat{\location} = \codequote{x~\texttt{=}~\tload^{\tau}~e}$}.
Then $\safeet(\tree(s, [t]))$ if there exist $\mo$, $\symlocalstore'$, and $t$ such that:\footnote{
The definition of well-formedness for symbolic memory mirrors~\Cref{def:well-formed}.
}
(1) $\mo \in \symmemory$,
(2) $\symmemory$ is well-formed,
(3) \textup{$\evalexprty{e}{\symlocalstore}{\bvsort{64}} = p$},
(4) $\unsat{\pc \land \neg \rangecond(p, |\tau|, \mo)}$,
(5) $\symlocalstore[x \mapsto \symmoread{\tau}{\mo}{p}] \equiv \symlocalstore'$,
(6) $t.\attrstate = \symstate{\nextinstr{\location}}{\symlocalstore'}{\symstack}{\symmemory}{\pc}$, and
(7) $\safeet(t)$.
\end{interfacelemma}
